1. En un grupo de 100 personas el 84\% tiene los ojos castaños, 
el 24\% tiene los cabellos castaños y el 15\% tiene los ojos y los cabellos castaños. 
¿Cuál es la probabilidad de que una persona elegida al azar no tenga ni los ojos ni 
los cabellos castaños?

Utilizaremos el \textbf{Principio de Inclusión-Exclusión} para calcular la probabilidad 
de que una persona tenga al menos una de las características.

Sea A el evento de que una persona tenga los ojos castaños.
$$P(A) = 0.84$$

Sea B el evento de que una persona tenga los cabellos castaños.
$$P(B) = 0.24$$

Sea $A\cap B$ el evento de que una persona tenga ambos 
(ojos y cabellos castaños).

$$P(A \cap B) = 0.15$$


Primero, calculamos la probabilidad de que una persona tenga ojos castaños o 
cabellos castaños:

$$P(A \cup B) = P(A) + P(B) - P(A \cap B)$$
$$P(A \cup B) = 0.84 + 0.24 - 0.15$$
$$P(A \cup B) = 1.08 - 0.15$$
$$P(A \cup B) = 0.93$$

Esta es la probabilidad de que una persona tenga al menos una de las características. 
Para encontrar la probabilidad de que una persona no tenga ni los ojos ni los 
cabellos castaños, simplemente restamos la probabilidad total menos la probabilidad 
de que tenga al menos una característica:

$$P(\text{Ni A ni B}) = 1 - P(A \cup B)$$
$$P(\text{Ni A ni B}) = 1 - 0.93$$
$$P(\text{Ni A ni B}) = 0.07$$

Por lo tanto, la probabilidad de que una persona elegida al azar no tenga ni los 
ojos ni los cabellos castaños es del 0.07.

\newpage

2. Dos dados se lanzan al azar (considera que son distinguibles y convencionales) 
¿Cuál es la probabilidad de que la suma de sus caras sea menor a 11?

Cuando se lanzan dos dados distinguibles, cada dado tiene 6 caras. 
El número total de resultados posibles es el producto del número de 
resultados para cada dado.

Número total de resultados posibles: 36

Es más sencillo calcular la probabilidad del evento complementario que la suma de las 
caras sea 11 o más, y luego restar esta probabilidad de 1.

Los resultados donde la suma es 11 o más son:

Suma = 11: $(5, 6), (6, 5)$

Suma = 12: $(6, 6)$

Número de resultados donde la suma es 11 o más: 3 resultados.

Sea $S$ el evento de que la suma de las caras de los dos dados sea 11 o más.

Ahora, calculamos la probabilidad de que la suma sea 11 o más:
$$P(S) = \frac{\text{Número de resultados con suma} \geq 11}{\text{Número total de resultados}}$$
$$P(S) = \frac{3}{36}$$
$$P(S) = \frac{1}{12}$$

Finalmente, calculamos la probabilidad de que la suma sea menor a 11 ($\lnot  S$):
$$P(\lnot S) = 1 - P(S)$$
$$P(\lnot S) = 1 - \frac{1}{12}$$
$$P(\lnot S) = \frac{11}{12}$$

Por lo tanto, la probabilidad de que la suma de las caras de dos dados distinguibles 
sea menor a 11 es de 11/12.

\newpage

3. Elegimos una carta al azar en una baraja de naipes (con 52 cartas). 
¿Cuál es la probabilidad de que la carta sea roja o tenga una figura?

Este problema también utiliza el \textbf{Principio de Inclusión-Exclusión}.

Número total de cartas en la baraja: 52

Sea \textbf{R} el evento de que la carta elegida sea \textbf{roja}.

Hay 2 palos rojos (corazones y diamantes), y cada palo tiene 13 cartas.

Número de cartas rojas: $2 \times 13 = 26$. Entonces
$$P(R) = \frac{26}{52} = \frac{1}{2}$$

Sea \textbf{F} el evento de que la carta elegida sea una figura 
(K, Q, J). Hay 3 figuras por cada uno de los 4 palos.

Número de figuras: $3 \times 4 = \textbf{12}$
$$P(F) = \frac{12}{52} = \frac{3}{13}$$

Sea $R \cap F$ el evento de que la carta sea roja y figura. Hay 3 figuras en 
corazones y otras 3 figuras en diamantes.

Número de figuras rojas: $3 + 3 = 6$
$$P(R \cap F) = \frac{6}{52} = \frac{3}{26}$$

Ahora, aplicamos el Principio de Inclusión-Exclusión para encontrar la probabilidad de 
que la carta sea roja o una figura:

$$P(R \cup F) = P(R) + P(F) - P(R \cap F)$$
$$P(R \cup F) = \frac{26}{52} + \frac{12}{52} - \frac{6}{52}$$
$$P(R \cup F) = \frac{26 + 12 - 6}{52}$$
$$P(R \cup F) = \frac{32}{52}$$
$$P(R \cup F) = \frac{8}{13}$$

Por lo tanto, la probabilidad de que la carta sea roja o tenga una figura es de 8/13.

\newpage

4. De forma simultánea elegimos dos cartas al azar en una baraja de naipes 
(con 52 cartas). ¿Cuál es la probabilidad de que formen un par?

Para este problema, usaremos combinaciones porque el orden en que se eligen 
las cartas no importa.

Número total de formas de elegir 2 cartas de 52:

Usamos la fórmula de combinaciones $C(n, k) = \frac{n!}{k! \cdot (n-k)!}$

$C(52, 2) = \frac{52!}{2! \cdot (52-2)!} = \frac{52 \times 51}{2} = 26 \times 51 = 1326$.

Este es el denominador de nuestra probabilidad.

Número de formas de elegir un par:
    
Para formar un par, primero debemos elegir el valor del par. Hay 13 
valores posibles (As, 2, 3, ..., 10, J, Q, K).

$C(13, 1) = 13$ formas de elegir el valor.

Una vez elegido el valor (por ejemplo, "ases"), necesitamos elegir 2 cartas de ese valor de las 4 
disponibles (por ejemplo, los 4 ases).

$C(4, 2) = \frac{4!}{2! \cdot (4-2)!} = \frac{4 \times 3}{2} = 6$ 
formas de elegir 2 cartas de un valor específico.

Número total de pares posibles:
    
(Número de formas de elegir el valor) $\times$ (Número de formas de elegir 2 cartas de ese valor)
$13 \times 6 = 78$.
Este es el numerador de nuestra probabilidad.
    
Finalmente, calculamos la probabilidad de formar un par:
$$P(\text{Par}) = \frac{\text{Número de formas de formar un par}}{\text{Número total de formas de elegir 2 cartas}}$$
$$P(\text{Par}) = \frac{78}{1326}$$

Para simplificar la fracción, podemos dividir por 6:
$$P(\text{Par}) = \frac{\frac{78}{6}}{\frac{1326}{6}} = \frac{13}{221}$$

Por lo tanto, la probabilidad de que las dos cartas elegidas formen un par es de 13/221.

\newpage

5: Si dos personas tienen que sus cumpleaños ocurren en un día al azar 
(considera años de 365 días), ¿Cuál es la probabilidad de que las dos personas tengan el mismo cumpleaños?

En este problema, consideramos 365 días posibles para el cumpleaños de cada persona.

Número total de combinaciones posibles para los cumpleaños de dos personas:

La primera persona puede tener su cumpleaños en cualquiera de los 365 días.

La segunda persona también puede tener su cumpleaños en cualquiera de los 365 días.

Por lo tanto, el número total de pares de cumpleaños posibles es: $365 \times 365 = 133,225$.
Este es el denominador de nuestra probabilidad.

Número de formas en que las dos personas tengan el mismo cumpleaños:

Si la primera persona nace en un día específico (ej. 1 de enero), la segunda persona debe nacer en ese 
mismo día (1 de enero) para que tengan el mismo cumpleaños.

La primera persona puede nacer en cualquiera de los 365 días.

Para que la segunda persona tenga el mismo cumpleaños, solo hay una opción para su día de nacimiento 
(el mismo que la primera).

Así, el número de formas en que ambos tengan el mismo cumpleaños es: 
        
$365 \text{ (opciones primer cumpleaños)} \times 1 \text{(opción para el segundo)} = 365$.

Este es el numerador de nuestra probabilidad.

Finalmente, calculamos la probabilidad de que las dos personas tengan el mismo cumpleaños:
$$P(\text{Mismo Cumpleaños}) = \frac{\text{Número de formas de tener el mismo cumpleaños}}{\text{Número total de combinaciones de cumpleaños}}$$
$$P(\text{Mismo Cumpleaños}) = \frac{365}{365 \times 365}$$
$$P(\text{Mismo Cumpleaños}) = \frac{1}{365}$$

Por lo tanto, la probabilidad de que dos personas tengan el mismo cumpleaños es de 1/365.

\newpage